\section{Metodología}
La metodología aplicada se fundamentó en el marco ágil Scrum, reconocido por su flexibilidad y orientación a la entrega continua (Alarcón et al., 2021). El trabajo se organizó en sprints de dos semanas, donde cada iteración culminaba con un incremento funcional del sistema. Se llevaron a cabo reuniones de planificación, daily scrums y revisiones de sprint, garantizando retroalimentación constante.

Desde el punto de vista técnico, el backend se desarrolló en Laravel, seleccionado por su arquitectura MVC y sus herramientas nativas como Eloquent ORM, Blade y middleware de seguridad (Subecz, 2021). La elección se sustentó también en estudios comparativos que muestran la accesibilidad de Laravel frente a Symfony en equipos pequeños (Alrahhal & Alzoubi, 2021). En el frontend, React fue adoptado gracias a su enfoque en componentes reutilizables y su curva de aprendizaje intermedia (Kaur & Kaur, 2020).

La aplicación móvil, inicialmente desarrollada para Android, integra funcionalidades clave como notificaciones push y carga de evidencias multimedia. Estudios previos demuestran que React Native es una solución idónea para expandir la compatibilidad a iOS sin necesidad de reescribir la lógica principal (Silva et al., 2022).

Para garantizar calidad y seguridad, se implementaron procesos de pruebas automatizadas en componentes críticos, siguiendo buenas prácticas de integración continua y despliegue seguro (Redalyc, 2020).